\chapter{The Jacobian}
\label{mk-ch-jacobian}

\begin{definition}[Jacobian]
  \label{mk-def-jacobian}
  \dcref{III.6,custom}
  \uses{mk-thm-pic0-abelian}
  \cite{milne-av,kleiman-picard}
  The Jacobian of \(C/k\) is
  \[
    \operatorname{Jac}(C):=\operatorname{Pic}^0_{C/k}.
  \]
\end{definition}

\begin{remark}
  The definition is independent of a chosen rational point.  A point is used
  only to write the Abel--Jacobi map explicitly; changing it translates that
  map.
\end{remark}

\begin{theorem}[Jacobian package]
  \label{mk-thm-jacobian-package}
  \dcref{III.6,5.13,5.14,custom}
  \uses{mk-def-jacobian,mk-thm-abel-jacobi-universal,mk-thm-pic0-abelian}
  \cite{milne-av,kleiman-picard}
  By \ref{mk-thm-pic0-abelian}, the scheme \(\operatorname{Jac}(C)\) is an
  abelian variety of dimension \(g(C)\).  The map \(\iota_{P_0}\) is independent
  of the chosen base point up to translation, and it satisfies the universal
  property in \ref{mk-thm-abel-jacobi-universal}.
\end{theorem}

\begin{proof}
  The Picard-variety theorem identifies \(\operatorname{Pic}^0_{C/k}\) with
  an abelian variety and computes its dimension as \(g(C)\).  The
  Abel--Jacobi universal property then gives the stated factorisation for
  every pointed map from \(C\) to an abelian variety.  If \(P_1\) replaces
  \(P_0\), the two maps differ by the translation
  \(Q\mapsto Q-[\mathcal O_C(P_1-P_0)]\), by the tensor law.
\end{proof}

\begin{theorem}[Base change]
  \label{mk-thm-jacobian-base-change}
  \dcref{III.6,custom}
  \uses{mk-def-jacobian,mk-thm-picard-representability}
  For every field extension \(k\subset k'\),
  \[
    \operatorname{Jac}(C)\times_k k'\cong
    \operatorname{Jac}(C\times_k k').
  \]
  The isomorphism is canonical and compatible with towers of extensions and
  with Abel--Jacobi maps.
\end{theorem}

\begin{proof}
  Picard representability and the cohomology-and-base-change theorem identify
  the base change of each degree component with the corresponding component
  for \(C_{k'}\).  The identity section and tensor law are preserved by this
  identification, so the identity components agree.  Restricting to degree
  zero gives the canonical isomorphism of Jacobians, and the construction of
  the diagonal line bundle shows compatibility with Abel--Jacobi maps.
\end{proof}
